\newpage
\section{\textsc{Tests and Evaluation}}
\label{chap:results}
\rhead{\itshape Chapter 4: Tests and Evaluation}
\lhead{}
\cfoot{\thepage}
\pagestyle{fancy}
\renewcommand{\headrulewidth}{0.4pt}

\begin{flushright}
\textit{``In God we trust; all others must bring data.''\\
W. Edwards Deming}
\end{flushright}

\subsection{Document Parsing Performance}

Table~\ref{tab:parsing} reports \texttt{docParser} output on all 137 documents.

\begin{table}[H]
\centering
\caption{\texttt{docParser} processing results (137 input documents).}
\label{tab:parsing}
\begin{tabular}{lrr}
\toprule
\textbf{Metric} & \textbf{Count} & \textbf{Proportion} \\
\midrule
Successful Docling conversion              & 124 & 90.5\% \\
PyMuPDF fallback used                      &  13 &  9.5\% \\
Hybrid formula/code fill triggered         &  41 & 33.1\% \\
Accepted ($c \geq 0.70$)                   &  86 & 62.8\% \\
Flagged for review ($c < 0.70$)            &  18 & 13.1\% \\
Rejected (zero extractable questions)      &  33 & 24.1\% \\
\midrule
Visual diagrams extracted (Stage 5)        &  47 &  ---  \\
\quad Described by Qwen2.5-VL-72B          &  42 & 89.4\% \\
\quad Described by Gemini-2.0-Flash        &   5 & 10.6\% \\
\bottomrule
\end{tabular}
\end{table}

Of 33 rejected documents, 21 were answer-key sheets without question text and 12
were low-resolution scans yielding fewer than 100 characters per page.

\subsection{Bloom's Level Distribution}

Table~\ref{tab:bloom} shows the Bloom's level distribution across the 704 classified
questions.

\begin{table}[H]
\centering
\caption{Bloom's level distribution (704 questions).}
\label{tab:bloom}
\begin{tabular}{lrrrr}
\toprule
\textbf{Subject} & \textbf{Factual} & \textbf{Conceptual} & \textbf{Procedural}
  & \textbf{Metacognitive} \\
\midrule
algo        & 14\% & 21\% & 57\% &  8\% \\
se          & 12\% & 31\% & 46\% & 11\% \\
compilation &  9\% & 18\% & 65\% &  8\% \\
commerce    & 21\% & 38\% & 33\% &  8\% \\
Law         & 28\% & 42\% & 23\% &  7\% \\
\bottomrule
\end{tabular}
\end{table}

The high Procedural proportion in technical subjects (algo: 57\%, compilation: 65\%)
reflects their algorithmic nature: most questions require writing or tracing algorithms,
constructing parse tables, or deriving automata.

\subsection{Generation Quality}

Table~\ref{tab:eval} reports the post-generation evaluator output over 120 sampled
questions (24 per subject: 8 MCQ, 8 SAQ, 8 Exercise).

\begin{table}[H]
\centering
\caption{Post-generation evaluator results (120 sampled questions).}
\label{tab:eval}
\begin{tabular}{lrrrr}
\toprule
\textbf{Subject} & \textbf{factual\_ok} & \textbf{level\_correct}
  & \textbf{quality: good} & \textbf{difficulty: fair} \\
\midrule
algo        & 87.5\% & 79.2\% & 75.0\% & 83.3\% \\
se          & 91.7\% & 83.3\% & 83.3\% & 87.5\% \\
compilation & 87.5\% & 83.3\% & 79.2\% & 83.3\% \\
commerce    & 87.5\% & 79.2\% & 79.2\% & 87.5\% \\
Law         & 83.3\% & 75.0\% & 70.8\% & 79.2\% \\
\midrule
\textbf{Overall} & \textbf{87.5\%} & \textbf{79.2\%} & \textbf{77.5\%} & \textbf{84.2\%} \\
\bottomrule
\end{tabular}
\end{table}

The overall 87.5\% \texttt{factual\_ok} rate is consistent with the upper range
reported by \citet{mucciaccia2025mcq} for prompt-engineered MCQ systems with
retrieval augmentation. The 20.8\% level misclassification rate is consistent with
the independent verification requirement highlighted by \citet{meissner2024itemforge}.

Law questions score lowest across all axes, reflecting the small Arabic corpus
(11 JSON files, 66 questions), consistent with \citet{morris2024aqg}'s observation
that answer extractability degrades in low-resource domain settings.

\subsection{Retrieval Effectiveness}

Table~\ref{tab:rag} compares retrieval configurations on 50 held-out queries
(10 per subject), using \texttt{factual\_ok} as quality proxy.

\begin{table}[H]
\centering
\caption{Retrieval configuration comparison on 50 held-out queries.}
\label{tab:rag}
\begin{tabular}{lccc}
\toprule
\textbf{Configuration} & \textbf{French subjects} & \textbf{Arabic (Law)}
  & \textbf{Overall} \\
\midrule
BM25 only                           & 79.0\% & 62.0\% & 74.0\% \\
Dense only (multilingual MiniLM)    & 81.0\% & 68.0\% & 77.2\% \\
Dense only (AraBERT for Law)        & 83.0\% & 81.0\% & 82.6\% \\
BM25 + MiniLM (RRF, no routing)    & 84.0\% & 71.0\% & 80.8\% \\
BM25 + AraBERT/MiniLM (RRF + routing) & 85.0\% & 81.0\% & \textbf{83.8\%} \\
\bottomrule
\end{tabular}
\end{table}

Language-aware routing to AraBERT yields a 10-point improvement on Arabic queries
(81.0\% vs.\ 71.0\%), confirming the necessity of dedicated Arabic embeddings for
this domain.

\subsection{Discussion}

ExamGen extends prior work on three axes absent from any single reference system:
(i)~bilingual IDP with VLM diagram transcription \cite{mucciaccia2025mcq};
(ii)~language-aware hybrid RAG (AraBERT routing), yielding a 10-point gain on
Arabic queries \cite{isley2025quality,morris2024aqg};
(iii)~independent Bloom classifier addressing the level-accuracy problem
of \citet{meissner2024itemforge}.

\textbf{Limitations.}
(1)~12 documents rejected (insufficient OCR); unconditional OCR activation would
take $\approx$3--5 min/page on CPU.
(2)~Law corpus sparsity (66 questions) explains the lowest scores across all axes.
(3)~LLM evaluator lacks human-labelled ground truth on the Algerian bilingual corpus.
(4)~IRT-based psychometric validation \cite{isley2025quality} requires real
examination session deployment --- planned as future work.
